
%\setcounter{demsode}{1} %%% THỨ TỰ ĐỀ THẦY - CÔ ĐÃ NHẬN
%
%\begin{name}
%	{Biên soạn: Thầy Nguyễn Cường\\
%	Phản biện: Thầy PT Sinh}
%	{ÔN TẬP HỌC KỲ II KHỐI 10 NĂM 2024}
%\end{name}

%\Opensolutionfile{ans}[ans/0-CTST-TN-SO-1]
%\setcounter{ex}{0}
\caulc
%%% CÂU 1
\begin{ex} %[DA-HK2-K10,11-2024 - Nguyễn Cường]%[0D8N2-4]
	Có bao nhiêu cách chọn $3$ học sinh từ một nhóm gồm $8$ học sinh?
	\choice
	{$8^3$}
	{\True $\mathrm{C}_8^3$}
	{$\mathrm{A}_8^3$}
	{$3^8$}
	\loigiai{
	
	}
\end{ex}

%%% CÂU 2
\begin{ex}%[DA-HK2-K10,11-2024 - Nguyễn Cường]%[0D8H1-5]
	Có bao nhiêu số tự nhiên có bốn chữ số?
	\choice
	{$5040$}
	{$4536$}
	{$10000$}
	{\True $9000$}
	\loigiai{
	
	}
\end{ex}
%%% CÂU 3
\begin{ex} %[DA-HK2-K10,11-2024 - Nguyễn Cường]%[0D7H3-2]
	Số nghiệm của phương trình $\sqrt{x^2+3}=3x-1$ là
	\choice
	{$0$}
	{\True $1$}
	{$2$}
	{$3$}
	\loigiai{
		
	}
\end{ex}

%%% CÂU 4
\begin{ex}%[DA-HK2-K10,11-2024 - Nguyễn Cường]%[0D8N1-1]
	Bạn Vy có $3$ cây viết chì, $8$ cây viết bi xanh và $2$ cây viết bi đỏ trong hộp bút, các cây viết phân biệt. Có bao nhiêu cách để bạn Vy chọn ra một cây viết?
	\choice
	{\True $13$}
	{$11$}
	{$48$}
	{$10$}
	\loigiai{
		
	}
\end{ex}
%%% CÂU 5
\begin{ex} %[DA-HK2-K10,11-2024 - Nguyễn Cường]%[0H9H3-2]
	Trong mặt phẳng $Oxy$, viết phương trình đường thẳng đi qua điểm $A(3;2)$ và nhận $\overrightarrow{n}=(2;-4)$ làm véc-tơ pháp tuyến.
	\choice
	{$2x+y-8=0$}
	{\True $x-2y+1=0$}
	{$x-2y-7=0$}
	{$3x-2y+4=0$}
	\loigiai{
	
	}
\end{ex}

%%% CÂU 6
\begin{ex}%[DA-HK2-K10,11-2024 - Nguyễn Cường]%[0H9N5-2]
	Trong mặt phẳng $Oxy$, phương trình nào sau đây là phương trình chính tắc của một Elip?
	\choice
	{$\dfrac{x}{9}+\dfrac{y}{8}=1$}
	{\True $\dfrac{x^2}{9}+\dfrac{y^2}{1}=1$}
	{$\dfrac{x^2}{2}+\dfrac{y^2}{3}=1$}
	{$\dfrac{x^2}{9}-\dfrac{y^2}{8}=1$}
	\loigiai{
	
	}
\end{ex}
%%% CÂU 7
\begin{ex} %[DA-HK2-K10,11-2024 - Nguyễn Cường]%[0H9N3-1]
	Trong mặt phẳng $Oxy$, đường thẳng $-x+2y+2024=0$ có một véc-tơ pháp tuyến là véc-tơ nào sau đây?
	\choice
	{$\overrightarrow{n}=(1;2)$}
	{$\overrightarrow{n}=(-1;-2)$}
	{\True $\overrightarrow{n}=(-1;2)$}
	{$\overrightarrow{n}=(-2;1)$}
	\loigiai{
	
	}
\end{ex}

%%% CÂU 8
\begin{ex}%[DA-HK2-K10,11-2024 - Nguyễn Cường]%[0D7H2-1]
	Tổng số nghiệm nguyên của bất phương trình $2x^2-5x+2\le 0$ là
	\choice
	{\True $3$}
	{$4$}
	{$5$}
	{$6$}
	\loigiai{
		
	}
\end{ex}

%%% CÂU 9
\begin{ex} %[DA-HK2-K10,11-2024 - Nguyễn Cường]%[0H9H4-4]
	Trong mặt phẳng $Oxy$, bán kính của đường tròn tâm $I(3;2)$ tiếp xúc với đường thẳng $\Delta\colon x+5y+1=0$ bằng
	\choice
	{$\sqrt{26}$}
	{$\dfrac{14\sqrt{26}}{13}$}
	{$5$}
	{\True $\dfrac{7\sqrt{26}}{13}$}
	\loigiai{
	
	}
\end{ex}

%%% CÂU 10
\begin{ex}%[DA-HK2-K10,11-2024 - Nguyễn Cường]%[0D0N1-3]
	Xét phép thử gieo một đồng xu cân đối và đồng chất ba lần. Xét biến cố $A\colon$\lq\lq Cả ba lần gieo cùng sấp hoặc cùng ngửa\rq\rq. Khi đó
	\choice
	{$n(A)=4$}
	{$n(A)=6$}
	{$n(A)=1$}
	{\True $n(A)=2$}
	\loigiai{
	
	}
\end{ex}
%%% CÂU 11
\begin{ex} %[DA-HK2-K10,11-2024 - Nguyễn Cường]%[0D7H1-2]
	Tam thức nào dưới đây luôn dương với mọi giá trị của $x$?
	\choice
	{$x^2-10x+2$}
	{$x^2-2x-10$}
	{\True $x^2-2x+10$}
	{$-x^2+2x+10$}
	\loigiai{
		
	}
\end{ex}

%%% CÂU 12
\begin{ex}%[DA-HK2-K10,11-2024 - Nguyễn Cường]%[0D0H2-2]
	Gieo một con xúc xắc cân đối và đồng chất hai lần. Tính xác suất để tích số chấm xuất hiện trên con xúc xắc trong hai lần giao là một số lẻ.
	\choice
	{$0{,}85$}
	{$0{,}5$}
	{\True $0{,}25$}
	{$0{,}75$}
	\loigiai{
	
	}
\end{ex}
\Closesolutionfile{ans}
\indapan{6}{ans/0-CTST-TN-SO-1}

\cauds
%\Opensolutionfile{ans}[ans/0-CTST-DS-SO-1]
%\setcounter{ex}{0}

%%% CÂU 1
\begin{ex} %[DA-HK2-K10,11-2024 - Nguyễn Cường]%[0D7V2-3]
	Cho biểu thức $f(x)=\dfrac{x-3}{x^2+7x+6}$. Xét tính đúng sai của các mệnh đề sau:
	\choiceTF
	{\True Tập xác định của hàm số là $\mathscr{D}=\mathbb{R}\setminus \{-1;-6\}$}
	{$f(x)=0\Leftrightarrow\hoac{&x=1\\&x=-6}$}
	{Với $x\in (-\infty;-6)\cup (-1;3)$ thì $f(x)>0$}
	{Với $x\in (-6;-1)$ thì $g(x)=f(x)-1>0$}
	\loigiai{
		
	}
\end{ex}

%%% CÂU 2
\begin{ex}%[DA-HK2-K10,11-2024 - Nguyễn Cường]%[0D8V2-7]
	Một đoàn tàu nhỏ có $3$ toa khách đỗ ở sân ga. Có $3$ hành khách không quen biết cùng bước lên tàu. Xét tính đúng của các mệnh đề sau:
	\choiceTF
	{Số khả năng khách lên tàu tùy ý là $9$ khả năng}
	{Số khả năng $3$ hành khách lên cùng một toa là $1$ khả năng}
	{\True Số khả năng mỗi khách lên một toa là $6$ khả năng}
	{\True Số khả năng có $2$ hành khách lên cùng một toa, hành khách thứ ba thì lên toa khác là $18$}
	\loigiai{
		\begin{itemchoice}
			\itemch Mỗi hành khách có $3$ cách lên tàu tùy ý nên có $3^3=27$ cách.
			\itemch Cả $3$ hành khách cùng lên một toa tàu có $3$ cách vì có $3$ toa tàu khác nhau.
			\itemch Số khả năng mỗi khách lên một toa là $6$ khả năng là số hoán vị của $3$ phần tử nên có $3!=6$ cách.
			\itemch Chọn $2$ hành khách trong số $3$ hành khách có $\mathrm{C}_3^2=3$ cách.\\
			Ứng với mỗi cách chọn đó có $3$ cách chọn toa tàu cho hành khách đi lên.\\
			Hành khách cuối cùng có $2$ cách để lên tòa tàu còn lại.\\
			Do đó, có $3\cdot 3\cdot 2=18$ cách.
		\end{itemchoice}
	}
\end{ex}

%%% CÂU 3
\begin{ex} %[DA-HK2-K10,11-2024 - Nguyễn Cường]%[0D8H2-7]
	Lớp 10A1 trường THPT Ten Lơ Man có $44$ học sinh, trong đó có $18$ học sinh sinh nữ và $26$ học sinh nam.
	\choiceTF 
	{Số cách sắp xếp $18$ học sinh sinh nữ và $26$ học sinh nam thành một hàng dọc là $18!\cdot 26!$}
	{Có $4!$ cách xếp khác nhau cho $4$ học sinh ngồi vào $6$ chỗ trên một bàn dài}
	{GVCN cần lập ban cán sự lớp gồm $1$ lớp trưởng, $1$ lớp phó học tập và $1$ lớp phó kỉ luật. Có $3!$ cách lập}
	{\True Lớp trưởng cần chọn ra $5$ bạn và xếp thứ tự thuyết trình cho $5$ bạn này. Có $\mathrm{A}_{44}^5$ cách chọn}
	\loigiai 
	{
	
	}
\end{ex}

%%% CÂU 4
\begin{ex}%[DA-HK2-K10,11-2024 - Nguyễn Cường]%[0H9H3-1]%[0H9H3-2]%[0H9H3-3]%[0H9H3-4]
	Trong mặt phẳng $Oxy$, cho hai đường thẳng $d_1\colon 3x+4y+1=0$ và $d_2\colon\heva{&x=15+12t\\&y=1+5y}$. Xét tính đúng sai của các mệnh đề sau:
	\choiceTF
	{\True Đường thẳng $d_1$ có một véc-tơ pháp tuyến là $\overrightarrow{n}=(3;4)$}
	{\True Phương trình tổng quát của đường thẳng $d_2\colon 5x-12y-63=0$}
	{Góc tạo bởi hai đường thẳng $d_1$ và $d_2$ có số đo bằng $60^\circ$}
	{\True Giao điểm của hai đường thẳng $d_1$ và $d_2$ là $A\left(\dfrac{30}{7};\dfrac{-97}{28}\right)$}
	\loigiai{
		
	}
\end{ex}
%\Closesolutionfile{ans}
%\indapan{4}{ans/0-CTST-DS-SO-1}

\caukq
%\Opensolutionfile{ans}[ans/0-CTST-KQ-SO-1]
%\setcounter{ex}{0}

%%% CÂU 1
\begin{ex}%[DA-HK2-K10,11-2024 - Nguyễn Cường]%[0D7H3-2]
	Tìm tổng các nghiệm của phương trình $\sqrt{-x^2+7 x+13}=5$.
	\shortans[]{$7$}
	\loigiai{
		
	}
\end{ex}

%%% CÂU 2
\begin{ex}%[DA-HK2-K10,11-2024 - Nguyễn Cường]%[0H9H3-8]
	Một trạm viễn thông $S$ có tọa độ $(5;1)$. Một người đang ngồi trên chiếc xe khách chạy trên đoạn cao tốc có dạng một đường thẳng $\Delta$ có phương trình $12x+5y-20=0$. Tính khoảng cách ngắn nhất giữa người đó và trạm viễn thông $S$. Biết rằng mỗi đơn vị độ dài tương ứng với $1$ km.
	
	\shortans[]{$3{,}46$}
	\loigiai{
	
	}
\end{ex}

%%% CÂU 3
\begin{ex}%[DA-HK2-K10,11-2024 - Nguyễn Cường]%[0D7H1-5]
Một vật chuyển động có vận tốc (m/s) được biểu diễn theo thời gian $t$ (s) bằng công thức $v(t)=\dfrac{1}{2}t^2-4t+10$. Trong $10$ giây đầu tiên, vận tốc của vật đạt giá trị nhỏ nhất bằng bao nhiêu?

\shortans[]{$2$}
\loigiai{

}	
\end{ex}

%%% CÂU 4
\begin{ex}%[DA-HK2-K10,11-2024 - Nguyễn Cường]%[0D7V2-7]
	Tổng chi phí $P$ (đơn vị nghìn đồng) để sản xuất $x$ sản phẩm được cho bởi công thức $P=x^2+30x+3300$; giá bán một sản phẩm là $170$ nghìn đồng. Số sản phẩm được sản xuất trong khoảng $(a;b)$ để đảm bảo nhà sản xuất có lãi (giả sử sản phẩm bán được hết). Tính $b-a$.
	
	\shortans[]{$80$}
	\loigiai{
	
	}	
\end{ex}

%%% CÂU 5
\begin{ex}%[DA-HK2-K10,11-2024 - Nguyễn Cường]%[0D0V2-5]
	Một hội đồng có đúng $1$ người là nữ. Nếu chọn ngẫu nhiên $2$ người từ hội đồng thì xác suất cả hai người đều là nam là $0{,}8$. Hội đồng có bao nhiêu người?
	\shortans[]{$10$}
	\loigiai{
		
	}	
\end{ex}

%%% CÂU 6
\begin{ex}%[DA-HK2-K10,11-2024 - Nguyễn Cường]%[0H9C5-6]
	Một tháp làm nguội của một nhà máy có mặt cắt là hình hypebol có phương trình $\dfrac{x^2}{30^2}-\dfrac{y^2}{50^2}=1$. Biết chiều cao của tháp là $120$ m và khoảng cách từ nóc tháp đến tâm đối xứng của hypebol bằng $\dfrac{1}{2}$ khoảng cách từ tâm đối xứng đến đáy. Tính bán kính đáy của tháp.
	
	\shortans[]{$38$}
	\loigiai{
	
	}	
\end{ex}

%\Closesolutionfile{ans}
%\indapan{6}{ans/0-CTST-KQ-SO-1}